\documentclass[10pt,a4paper]{article}
\usepackage[T1]{fontenc}
\usepackage[utf8]{inputenc}
\usepackage[margin=1in]{geometry}
\usepackage{amsmath,booktabs}
\usepackage[hidelinks]{hyperref}
\setlength{\parskip}{2pt}
\setlength{\emergencystretch}{2em}
\newcommand{\Host}{\mathrm{Host}}
\newcommand{\Recomp}{\mathrm{Recompute}}
\title{Beyond the Next Token: Native KV Recovery,\newline Re-preemption, and Peer Cost}
\author{Anonymous Author(s)}
\date{Research draft --- October 8, 2026}

\begin{document}
\maketitle
\begin{abstract}
A preempted language-model request can resume by loading available host KV or
recomputing its known history. We study this choice in one native serving
runtime while fixing victim selection, queues, admission, and execution
mechanisms. Two repeated single-event interventions show why decision-to-next-output
latency alone is insufficient evidence of service benefit. Host restoration is
locally faster at one matched recovery point; recomputation is faster at another
because it avoids re-preemption, but delays a peer and discards that peer's
intermediate work. A remaining-token rule produces exactly the same actions as
a length threshold and is removed. A subsequent KV-headroom rule improves one
request's completion while increasing mean completion time in both order
comparisons. Natural outputs and total work differ across actions, and separate
telemetry and CPU diagnostics leave substantial same-work timing drift
unexplained. These observations support a bounded mechanism account, not an
adaptive-policy improvement or a general action-ranking theorem.
\end{abstract}

\section{Problem and scope}
KV recovery is part of a continuing service trajectory. A request may have
already produced tokens, have only part of its history available in host memory,
and compete with other requests for its next compute allocation. The time to
load a cache or execute a prefill in isolation does not describe all these
consequences. Even a recovery that reaches its next output earlier can alter the
next victim, delay another request, or increase repeated computation.

We examine a deliberately narrow choice: for a request that the existing
scheduler has selected to resume, should it reuse ready host KV or rebuild the
same known history using native recomputation? The destination GPU, victim
rule, queue order, admission checks, token quantum, transfer concurrency, and
transfer chunking are fixed within every comparison. We do not choose which
request to evict, reorder recovery tasks, or change when new work is admitted.
Different later states and victims are possible consequences of the action,
even though the rules selecting them remain unchanged.

The empirical question is whether a local recovery improvement survives the
costs imposed on the complete request population. We do not assume a crossover,
an optimal online selector, or a state advantage over a fixed threshold.
Two initial natural-output interventions change exactly one
successful recovery allocation after matching the preceding observed history.
They expose both a locally faster Host action and a locally faster Recompute
action, with different effects on peers. The two points have different prefixes
and host coverage, so their rankings do not isolate the incremental information
in current load. A later fixed-output-count intervention on another card follows
an early Recompute lead through subsequent Host recovery and finds that the
lead disappears without changing any request's completion step.

This draft also retains two unsuccessful online candidates and two timing
diagnostics. The token-budget candidate has no action advantage over a length
rule. The headroom candidate changes an action and improves a particular
request, but does not establish a whole-service win. Same-policy repetitions
produce identical outputs and native schedules while taking substantially
different time. We report those differences without subtracting estimated
overheads or selecting favorable runs. The scope is a measured native-runtime
mechanism, not a submission-ready claim of a new winning serving policy.

\section{The native decision and its cost boundary}
\subsection{Known history, host coverage, and eligibility}
For request $r$ at decision time $t$, let
\[
 X_r(t)=P_r\Vert Y_r^{\leq t},\qquad K_r(t)=|X_r(t)|,
\]
where $P_r$ is the received input and $Y_r^{\leq t}$ contains only previously
generated tokens. Future output length and reference answers are unavailable to
the selector. Let $h_r$ be the prefix length reported ready by the native host
connector. Host loads that prefix and computes any missing suffix. Recompute
rebuilds the necessary KV from $X_r(t)$ in the same request object; it neither
resamples the old prefix nor emits recomputation tokens as new output.

With block size $b$, we conservatively require the capacity for the complete
current history before enabling the choice:
\begin{equation}
 \left\lceil K_r/b\right\rceil+B_{\rm reserved}+B_{\rm watermark}
 \leq B_{\rm free}.
 \label{eq:capacity}
\end{equation}
This supplements, rather than replaces, native allocation checks. A positive,
ready host hit and Eq.~\ref{eq:capacity} define the jointly eligible domain.
Partial host coverage is allowed: $h_r<K_r$ does not remove the need to rebuild
the missing suffix. A smaller initial chunk allocation is not evidence that
less complete KV is required for the next correct output.

The implementation uses the recorded vLLM 0.26.0 native CPU-offloading path.
The connector's \texttt{get\_num\_new\_matched\_tokens} performs the common lookup.
Host preserves its result. For Recompute, the thin adapter clears the matching
connector hit metadata and returns \texttt{(0, False)}, so the native scheduler
recomputes the original request. It does not introduce another executor or run
both branches. The existing \texttt{update\_state\_after\_alloc} callback marks
successful commitment; unsuccessful lookups are not counted as interventions.
Pending transfers, host misses, and failures of joint capacity retain their
native behavior and are recorded separately. Thus, a fixed Host policy can
still execute native recomputation on a host miss.

The adapter checks request-object identity, known tokens, sampling-parameter
object, generated-token count, status, and stop reason across selection.
It records a prefix hash, preserves monotonically extending output, and disables
selection during warmup. These are continuity checks, not elementwise KV,
logit, stochastic-RNG, or semantic-equivalence proofs. The measured workload
uses greedy decoding. Later output differences are reported rather than hidden
behind the state assertions.

\subsection{Measurements and sunk costs}
We record the time from a committed decision to the next strictly later output
of that request. This includes actual waiting, transfers, native computation,
and any intervening preemption. Multiple commits can share the same next
output; the analysis retains that association instead of treating them as
independent trials. We also inspect the following output gap and subsequent
completion, because one early token need not imply continuing progress.

For all externally arriving requests, including delayed admission, we report
\begin{equation}
 \overline F=\frac{1}{N}\sum_i(C_i-A_i),\qquad
 M=\max_i C_i-\min_i A_i,\qquad
 \Delta_{i,k}=O_{i,k}-O_{i,k-1}.
\end{equation}
Here $A_i$, $C_i$, and $O_{i,k}$ are external arrival, completion, and output
times. TTFT is separate from inter-output gaps. Output timestamps are observed
after synchronous \texttt{engine.step()} returns; they include driver and
instrumentation work and are not network-client measurements.

STORE performed before the choice is already paid. It remains in the complete
episode but is neither charged again to the incremental Host action nor counted
as a saving from Recompute. Subsequent STORE, LOAD, and repeated history work
remain in the measurements. CUDA transfer-event durations can overlap compute
and other transfers; their sums are diagnostics, not additive wall-clock costs.
Likewise, sums of gaps across concurrent requests are not elapsed service time.
Scheduled token positions measure native work, not GPU execution duration.

\section{Experimental method}
The completed main comparisons use one NVIDIA RTX PRO 6000 Blackwell Server
Edition, the existing OLMoE snapshot in BF16, PyTorch 2.11.0, Transformers
5.17.0, and vLLM 0.26.0. Execution is eager with the native Triton MoE backend,
FCFS, synchronous scheduling, context limit 4096, at most 256 requests, and a
2048-token step budget. Prefix caching and speculative decoding are disabled.
Native GPU profiling at utilization 0.9 produces 36,916 KV blocks of 16 tokens,
or 72.1015625 GiB; pinned host KV is 32 GiB. No explicit small KV limit is used
for these results.

Each group initializes one engine. Its arms use the same external arrivals,
input tokens, resource settings, full-population warmup, and drained cache reset.
The shared GPU lock serializes groups, and resource/queue reset assertions
check isolation between arms. Warmup uses native Host preference with the
selector disabled and generates to the caps; it is recorded separately from
the natural-EOS measurement. Initialization and lock waiting are outside episode
timing. End-to-end episode costs include service and transfer drain; the warmup
is not silently included as measured recovery work.

The mechanism workload has 256 requests: 64 frozen GSM8K questions and 192
full-document summary requests drawn cyclically from three complete documents
in the saved LongBench source~\cite{gsm8k,longbench}. Each submission has its
own cache identity. Repeated documents create a pressure probe, not 192
independent documents. Prompts are neither truncated nor padded or concatenated
to create pressure. Measured generation respects EOS, with short/long output
caps 256/1024 and no minimum generation length. All external arrivals are zero.

Normal-capacity QA trials at 64, 192, and 256 requests produced no recovery
choices. An independent MultiNews256 workload, using 27 complete news bundles,
also produced no preemption in any of its six H/R/L/L/R/H arms. Its observed
KV peak was 32,235/36,916 blocks (87.32\%). At the peak, 64 requests had already
retired; peak running population was only 192. Host residency reached 32 GiB
without any LOAD. These results show why a large sum of potential KV demands
or a full host cache does not establish an actionable recovery domain.

The summary workload does trigger both native paths. Its initial H/R/R/H
comparison had 46 jointly eligible commits, but one Recompute measurement
overlapped a known result-transfer operation. That arm remains marked as
confounded; no corrected latency is synthesized. Later groups were retrieved
only after completion and GPU release. Across-action output changes and timing
drift motivate the more specific interventions below.

\section{One action changes later execution opportunities}
\subsection{Intervention protocol}
A frozen target specifies external identity, preemption ordinal, known and
generated counts, and prefix hash. The intervention occurs at most once, only
after the naturally reached target passes joint eligibility and commits its
allocation. All other recoveries retain Host preference and the same native
fallback. A missing or mismatched target is recorded as inapplicable, not
manufactured by changing scheduling. Each reported arm reached its target once.
Within each four-arm group, normalized schedules, output histories, and prior
recovery histories match before the intervention. Normalization removes random
internal identifiers and times while preserving request order and token ranges.
All step and output indices below are zero-based.

\begin{table}[t]
\centering\small
\caption{Matched single-event interventions. Ranges contain the two repeated
observations, not confidence intervals. Peer delays share the target's decision
time origin; they must not be added to obtain system time.}
\label{tab:events}
\begin{tabular}{@{}llrr@{}}
\toprule
Target and known/host tokens & Observation & Host & Recompute \\
\midrule
E249: 3128/3056 & Target next output (ms) & 139.2--141.0 & 156.6--157.6 \\
 & E252 next output (ms) & 139.2--141.0 & 226.1--227.6 \\
E223: 3277/2656 & Target next output (ms) & 1017.8--1021.2 & 160.7--161.8 \\
 & E227 next output (ms) & 161.9--162.0 & 963.8--968.0 \\
\bottomrule
\end{tabular}
\end{table}

\subsection{E249: extra reconstruction delays another recovery}
The R/H/H/R group targets E249's first preemption, after 76 generated tokens,
at step 378. The common state has 282 free blocks, 196 required blocks, 192
running requests, seven waiting requests, and no pending LOAD. Host loads
3056 tokens and computes the remaining 72. Recompute executes 1856 positions
in step 378 and 1272 in step 379, rebuilding all 3128 known tokens. Both paths
produce the next token at step 379, but Host is 16.6--17.4 ms faster
(Table~\ref{tab:events}). The action replaces 382 MiB of target LOAD with 3056
additional native compute positions; this is work accounting, not a subtraction
of isolated costs.

At step 378, Recompute plus 192 single-token decoders exhausts the 2048-token
budget. Under Host, E252 can also start its native LOAD in that step, rebuild
its 15-token suffix in step 379, and output. Recompute delays those operations
one step, increasing E252's wait by 86.6--86.9 ms. The other 192 decoders are
still scheduled; the evidence is a displaced recovery opportunity, not an
inferred cancellation of their computation.

Same-action repetitions have identical complete schedules and outputs. Across
actions, 194 of 256 requests eventually differ in token content. The first
differences occur in requests that have never been preempted, while E249's first
65 resumed tokens agree. This is consistent with sensitivity to changed batch
computation, but without KV/logit measurements it identifies neither a numerical
root cause nor the absence of an implementation error. Recompute ultimately
adds 2039 output tokens and 15,434 repeated positions. Those episode differences
are not all direct costs of the one 3056-position substitution.

\subsection{E223: local progress transfers re-preemption to a peer}
The H/R/R/H group targets E223's second preemption at step 874, after 541
generated tokens. All arms have known length 3277, host hit 2656, 416 free
blocks, 205 required blocks, 161 running requests, and 24 waiting requests.
Host loads 2656 tokens, then computes 609 suffix positions. Twelve positions
remain when E223 is preempted again at step 876. A later LOAD and 13-position
suffix computation finally produce the next token at step 887. The two Host
commits refer to this same output.

Recompute instead receives 1880 and 1397 positions in steps 874--875 and
produces that token immediately, about 0.86 s earlier. The local improvement
contains avoided re-preemption and waiting, not merely a faster reconstruction
kernel. It also changes the peer E227's opportunities. E227 is a host miss and
uses the same native recomputation fallback in both arms. Under Host it rebuilds
3158 tokens in steps 874--875. Under Recompute it starts later, accumulates
2376 positions, is preempted at step 877, and must rebuild from zero again.
Its next output is delayed by about 0.80 s (Table~\ref{tab:events}).
The peer is re-preempted at step 877 while the finished-request sets still
match. Their first natural-stop difference occurs at step 884. Thus the earlier
re-preemption is not explained by that later stop; however, the peer rebuild
starts at 885, so final repeated-work totals still include later trajectory
and output-length changes.

E223 and E227 together execute 5031 additional positions in the Recompute
trajectory; the episode adds 2120 repeated positions and 814 output tokens after
other trajectory changes are included. Host's two target LOADs total 740 MiB,
but their approximately 14 ms CUDA-event sum cannot explain its one-second
next-output delay. All earlier STORE is retained. Cross-action content differs
for 143 requests despite identical target first resumed tokens. These
observations establish displaced work and re-preemption in this run, not an
equal-output service acceleration. Different history and host coverage also
prevent interpreting the opposite E249/E223 rankings as a controlled
load-dependent crossover.

\subsection{E246: an early lead disappears at the next Host recovery}
A later exploratory probe uses another GPU card, with 36,911 natively profiled
KV blocks, the frozen MultiNews320 inputs, and fixed generation counts of
256/1024 tokens instead of natural EOS. Two completed pairs execute H/R then
R/H; each arm completes 320 requests and emits 278528 tokens. The target was
selected after inspecting an earlier whole-policy fork, so this is action
isolation, not independent target selection. Each arm commits exactly one
target action, with every other recovery retaining Host preference and native
fallback. Within each pair, normalized schedules, output histories, and
recovery histories match before step 416. E246 then has 3165 known tokens,
including 125 generated tokens, a 1168-token Host hit, and 245 free blocks
against 198 required blocks.

Recompute reaches the next output in step 417 rather than Host's step 418,
leading by 60.671 and 65.264 ms in the two order comparisons. Host computes
1997 suffix positions; Recompute computes all 3165, adding 1168 positions.
Both trajectories resume E246 again at step 421 using Host, but their available
Host prefixes now differ: 3152 tokens after the initial Host action versus
1168 after Recompute. With known lengths 3167 and 3168, respectively, the
remaining native work is 15 versus 2000 positions, adding another 1985.
The episode therefore adds exactly $3153=1168+1985$ repeated positions, all
in E246, in both comparisons. Its first three resumed outputs advance, but
the 129th output (index 128) occurs in step 423 in both trajectories. At the
third recovery, step 426, both load 3168 tokens and compute two positions.
All 320 completion steps agree; all 319 peers also retain identical native
per-request schedules and output-emission steps. This is a loss of the target's
early progress advantage, not evidence of displaced peer computation.

Complete-service timing has no stable sign. Recompute minus Host makespan is
$+0.351$ s in H/R and $-0.471$ s in R/H; mean completion differences are
$+0.264$ and $-0.346$ s. Earlier STORE remains paid: every arm records
135499087872 STORE bytes, while Host and Recompute trajectories record
11634999296 and 11221860352 LOAD bytes. These byte counts and scheduled
positions are not additive wall-clock costs. Fixed output counts do not fix
contents or establish equal quality: 164 of 320 requests differ across actions
in each pair. No quality-adjusted service gain follows.

This closes the service-step-rule candidate without an adaptive-gain claim.
A completed passive field run, \path{oct08_e246_history53005_r01}, now follows
the native STORE path in both H/R arms; the pre-target checks pass and the
observer reports no errors. At step 416, Host allocation moves the STORE cursor
from chunk 197 to 73. Native jobs created in steps 417--418 contain five
missing chunks, in the half-open intervals $[73,74)$, $[189,192)$, and
$[194,195)$. Their native completion acknowledgements occur in steps 418 and
419, before the step-421 lookup returns a 3152-token Host hit. Recompute leaves
the cursor at 197; all four observed builder calls in steps 416--419 return no
STORE job, and its step-421 lookup returns 1168 tokens.

For this recovery interval, the field trace rules out newly created STORE jobs
merely remaining pending until the next lookup as the explanation for the
coverage difference. It does not establish continued residency, the cause of
the earlier missing chunks, or a general eviction mechanism. A subsequent single-entry causal probe, \path{oct08_e246_store_replay53005_r01},
executes the same preassigned Recompute action in both arms and changes only
native STORE-cursor replay (off/on). Both 320-request episodes finish, each
with 278528 output tokens and no failures, rejections, timeouts, or unfinished
requests. The pre-action histories match; the enabled arm performs exactly
one cursor rewind from 197 to 73 without changing the request state. Native
jobs in steps 416 and 417 store the same five missing chunks and are
acknowledged in steps 417 and 418. At the next lookup, step 421, both requests
have 3168 known tokens and 128 generated tokens, but Host coverage is 1168
without replay and 3152 with replay. Thus, in this event, the cursor change
is sufficient to restore the missing coverage; pending transfer alone is not
the cause.

Full-episode repeated positions decrease from 82376 to 80402. STORE bytes
remain 135499087872 in both arms; LOAD increases from 11221860352 to
11484004352 bytes. This is a change in when reusable state becomes available,
not a saving of previously paid STORE. The off/on mean completion times are
78.987693/79.021042 s and makespans are 121.744937/121.893368 s. The single
ordered pair does not establish a service improvement or a precise performance
effect: 36 requests already finish before intervention, with the on arm on
average 56.842 ms earlier. No timing correction is applied. Outputs differ
for 168 of 320 requests, despite equal output counts; quality equivalence is
not established. This supports a local implementation explanation, not an
inherent recomputation penalty or an independently supported online method.
Strong-baseline service value remains unestablished.

The enabled trace retains E246's progress advantage: its completion moves
from step 1911 to 1910, relative to both the off trace and the earlier Host
trace. Every peer retains its output-emission and completion steps; both
runs still require 2197 service steps. The long pause spans steps 545--1068
in both runs, moving from output indices 179--180 to 180--181. One output's
523-step advance therefore moves the pause rather than removing it. Some
peer prefill chunk boundaries change by one position; peer total work does
not. The enabled run still performs 1179 more repeated positions than the
earlier Host trace. This cross-group comparison concerns discrete work and
steps, not timing. On all 29 jointly eligible states in the enabled run,
both the frozen $q=1792$ step proxy and its supported current-budget version
recommend precisely the executed one Recompute and 28 Host actions. No
additional Recompute recommendation remains unexecuted on this trajectory.
This shadow agreement is not an online-policy experiment or an upper bound
on all possible policies. It removes the evidence-based reason to extend
this step-rule experiment at this operating point; it does not refute the
general recovery-choice problem.

\paragraph{Remaining opportunity and the cost of replenishment.}
A read-only follow-up of the latest fixed-Host baseline finds 23 Host-miss
fallbacks in 23 different requests, each the request's final recovery. Their
74932 repeated positions are paid reconstruction, not demonstrated savings
from future replenishment. An older natural-output baseline does contain two
Host misses for E243, with 2921 repeated positions at the second miss; its
other GPU and 256-request domain do not establish opportunity in the current
320-request fixed-output domain. Neither trace establishes that replenished
blocks would survive until reuse.

Replenishment is also not a cost-free correctness repair. In one constructed
CPU case with the saved native STORE builder/manager and official v0.26.0
LRU implementation, 15 Host blocks hold seven of A's eight chunks and all
eight of B's chunks. Rewinding A's cursor causes exactly one STORE, evicting
B's unprotected oldest chunk B0. Native consecutive lookup changes from
A=2/B=8 to A=8/B=0, although B1--B7 remain resident. The fixture supplies
compute progress and synchronous transfer completion; it measures no GPU,
copy time, later execution, or service gain. This counterexample limits the
interpretation of cursor replay to a cache re-admission policy with possible
peer costs. It is not a novelty claim or evidence that this eviction occurred
in the measured E246 episode. Both the trace extraction and the CPU case are
reproducible through \path{cpu_preparation/analyze_baseline_host_miss.py} and
\path{cpu_preparation/check_store_replay_eviction.py}.

\section{Online candidates: equivalence and a failed service improvement}
Fixed H and R select their respective action only in the jointly eligible
domain. The length rule L selects R when $K_r\leq1792$ and H otherwise. This
predeclared structural threshold is not a calibrated cost model and is not
claimed optimal.

The first state candidate B selected recomputation when the known history fit
within the observed remaining native token budget. In H/R/L/B/B/L/R/H, B and
L disagree on none of the 74 jointly eligible commits when replayed on the
observed states. Their four actual runs have identical 2024-step native
schedules, all 256 outputs, work, and transfer bytes. Relevant histories are
separated between 706 and at least 2736 tokens, while remaining budgets lie
between 1845 and 1884 on these L/B trajectories. B therefore adds no action
information in this domain. It and its budget observer were removed from the
active implementation; small timing differences are not an adaptive benefit.

The subsequent rule G adds one KV-headroom signal. For an eligible target, let
\begin{equation}
 m=B_{\rm free}-B_{\rm required}-B_{\rm reserved}-B_{\rm watermark}
       -g_{\rm background}-g_{\rm target}.
 \label{eq:margin}
\end{equation}
Here $g_{\rm background}$ counts current supported running requests at a block
boundary, and $g_{\rm target}$ is one when the target is at a boundary.
Support requires every running request to be one-token-ready, outside prefill
chunks, eligible for its native next decode step, and within the model limit.
Native computed counts advance at schedule end, so these snapshots describe
the allocations already made when the waiting request is inspected. G selects
R only when support holds, $K_r\leq1792$, and $m\geq0$; otherwise it selects H.
Unsupported observations retain eligibility but have a null margin. This is a
one-round estimate for the recorded set, not a promise against later admission,
release, or preemption.

\begin{table}[t]
\centering\small
\caption{The complete H/L/G/G/L/H comparison. Natural output work differs
across policies. Same-policy repeats preserve outputs and native schedules but
show timing drift; these are descriptive episode results.}
\label{tab:headroom}
\begin{tabular}{@{}lrrr@{}}
\toprule
Arm & Makespan (s) & Mean completion (s) & Output tokens \\
\midrule
H0 & 97.710 & 63.222 & 189436 \\
L1 & 106.434 & 64.547 & 189035 \\
G2 & 100.456 & 64.066 & 190178 \\
G3 & 114.256 & 65.269 & 190178 \\
L4 & 98.939 & 64.722 & 189035 \\
H5 & 111.679 & 63.615 & 189436 \\
\bottomrule
\end{tabular}
\end{table}

The six-arm group commits 60 eligible choices: 56 H and four R. Both G arms
commit eight H and one R; another 150 host-miss fallback commits across the
group remain native. For E252, G chooses H at known length 706 with margin
$-3$, then R at length 719 with margin $+31$. A usual monotone length threshold
cannot reproduce both actions. However, prefixes and host coverage also change;
this is action information, not identified incremental predictive value of state.

E252 finishes the same 56-token natural-stop output at 22.409/26.771 s under G,
versus 69.758/68.111 s under H. Its maximum gap falls from 47.613/42.057 s to
0.489/0.578 s. Yet the changed running order makes E249, rather than E252, the
native victim at step 382; E249's same-content output token at index 79 moves
from step 382 to 972. H also preempts E249 later, so these waits cannot be
subtracted to construct a system saving.
The victim changes at step 382, before the first completion-set difference
at 385 and the first differing terminal length at 389. The earlier completion
at 385 is E252 finishing the same full sequence. E249 itself has 3201 repeated
positions under G versus 3205 under H; this peer demonstrates displaced waiting,
not the source of G's aggregate repeated-work increase.

In both forward/reverse comparisons, G increases population mean completion
relative to H, by 0.844/1.654 s. Excluding E252, the other 255 requests average
1.033/1.823 s later completion. G executes 9999 additional positions: 742 more
output positions and 9257 more repeated-history positions. H/G have different
content for 196 requests. Table~\ref{tab:headroom} therefore provides no
whole-service win for G, and its natural-work differences and run-level drift
preclude an equal-work causal slowdown estimate. G is not retained as a winning
policy, and no threshold search follows this result.

\section{Timing diagnostics do not resolve same-work drift}
Within each of H, L, and G, repetitions above have identical schedules and
outputs, yet makespans vary by approximately 7.5--14 s. Differences concentrate
in small-batch tails, accompanied by driver-thread CPU time. CPU time can
include CUDA busy polling and is not pure host computation. Low-frequency GPU
samples show lower utilization in slow tails without lower sampled clocks;
they do not identify a cause. We performed two bounded diagnostic groups
without changing the recovery policy.

\begin{table}[t]
\centering\small
\caption{Fixed-Host timing controls with identical complete work and outputs.
The fixed tail is steps 1262--1992: 731 steps and 9026 scheduled positions.
Tail time is the sum of per-step wall-clock durations, not continuous-window elapsed time.
ON/OFF denotes only the GPU-query subprocess. The CPU group adds passive CPU
readings and keeps that subprocess ON.}
\label{tab:timing}
\begin{tabular}{@{}llrrr@{}}
\toprule
Group & Arm & Makespan (s) & Mean completion (s) & Tail step sum (s) \\
\midrule
GPU query & H0 ON  & 99.973 & 63.006 & 12.994 \\
 & H1 OFF & 107.917 & 64.673 & 21.161 \\
 & H2 OFF & 101.938 & 64.068 & 13.348 \\
 & H3 ON & 103.314 & 63.885 & 17.215 \\
\midrule
CPU readings & H0 & 95.818 & 62.159 & 10.662 \\
 & H1 & 110.880 & 63.955 & 24.427 \\
\bottomrule
\end{tabular}
\end{table}

The ON/OFF/OFF/ON experiment disables only the approximately 1 Hz
\texttt{nvidia-smi} querying subprocess. All four arms complete all requests,
with identical 1993-step schedules, outputs, and stopping metadata. Each has
14 eligible Host commits and 22 host-miss recomputation fallbacks. Switching
queries off increases the tail step-time sum by 8.168 s in the forward comparison
but reduces it by 3.867 s in the reverse comparison; the two OFF tails still
differ by 7.813 s.
The direction reversal does not support the query process as the established
cause, nor does it prove zero observer overhead. This control is stopped
without a frequency scan or additional repetitions. The sign reversal concerns
makespan and the tail step-time sum; mean completion is higher with queries OFF
in both comparisons.

The subsequent two-Host group adds passive CPU frequency, context-switch, and
container-throttling observations. It reproduces identical work but a 13.766 s
difference in the tail step-time sum and a 13.771 s driver-CPU difference.
Tail driver CPU totals are 10.673 and 24.444 s. The 11/25 tail frequency samples report the same last
CPU (60), approximately 3.2 GHz, and no observed CPU changes between valid
adjacent samples. Local cgroup throttling counters do not increase. These
readings do not distinguish the slow state: sampled sysfs frequency is not
effective execution frequency, polling gaps can hide migration, and ancestor
throttling is not excluded. Driver user time can still contain CUDA polling.
This CPU diagnostic is also stopped without affinity or frequency changes.

Across these controls each arm produces 189436 output tokens, schedules 860683
positions, repeats 72886 positions, and retains 103551074304 STORE bytes and
4349493248 LOAD bytes. No recovery-action difference explains the timing
variation. All raw times remain uncorrected. Neither diagnostic supports a
clean attribution of the existing between-policy makespan differences.

\section{Quality, related work, and remaining limits}
\subsection{Quality and applicability}
The frozen short-task rule extracts the last signed numeric value, normalizes
commas, compares it with the official GSM8K answer, and additionally requires
natural stop for the primary score. It is a defined numeric metric, not a
reasoning or factuality certification. The single-event and headroom groups
score 26/64 under this rule; equal scores do not imply equal outputs. In the G
group four short outputs change, while the correct-identity set remains the
same. The CPU-control outputs match the previously scored Host outputs exactly;
their inherited score is not new quality evidence.

The initial three-document summary audit found unsupported details, source
contradictions, repetition, and incomplete endings among six fixed identities.
Natural stop is therefore not synonymous with useful summary service. Those
findings are not a population error-rate estimate or a ranking of policies.
For independent MultiNews256, the official LongBench ROUGE-L behavior produces
a 27-bundle macro score of 27.243 and a request mean of 27.247, identically in
all six arms; four of 192 summaries reach the cap. Fixed review of three
identities still finds factual errors and limitations in the supplied references.
Since that workload has no recovery actions, neither equal outputs nor scores
demonstrate correctness or benefits of an exercised recovery mechanism.

A frozen MultiNews320 successor adds 64 complete long requests to the unchanged
256-request prefix, with the same 27 news bundles, EOS behavior, and caps.
It increases the allowed sequence count to 320, retaining native GPU profiling
and the 2048-token budget. Its H/R/R/H probe completes all 1280 requests with one jointly eligible
recovery per arm and no fallback. Native profiling yields 36,894 KV blocks.
At step 399, E307 resumes with 2875 known tokens and a 2864-token host hit;
observed histories and schedules match before that step. Both paths output
the same next token in step 400. Host times are 167.036/275.793 ms and
Recompute times are 179.388/178.821 ms: even the local ordering reverses.
Makespans in execution order are 69.432/60.201/68.320/58.088 s, and mean
completion times are 31.165/34.752/31.826/33.555 s. Thus, the between-policy
ordering of makespan, mean completion, and mean TTFT also reverses.

Same-policy repetitions have identical complete outputs, schedules, and work.
All 106 requests completed before the action are slower in the second arm of
each pair, irrespective of recovery policy. Recompute changes 167 outputs and
adds 31 output tokens plus 2864 repeated-history positions relative to Host.
Bundle-macro ROUGE-L is 27.182661/27.182155 for H/R, and all short-task scores
are 27/64. Close scores do not establish equivalent quality. These four arms
provide no stable service gain; the bounded candidate is closed without more
repetitions or threshold tuning.

The slow Host recovery has a 195.215 ms decision step, versus 83.303 ms in its
repeat, despite approximately 6.97 ms LOAD CUDA-event duration in both. Existing
trace timestamps locate the difference after the scheduler observation, not
within the recorded lookup or scheduling interval. A no-recovery control at
step 266 also differs by 125.540 ms, although its work shape differs from
step 399. The fluctuation is therefore not exclusive to recovery.

One subsequent fixed-Host diagnostic records the native transfer-start call
in all 1437 formal steps. The recovery call takes 0.435 ms, the decision step
70.765 ms, and the next output arrives after 139.839 ms; the earlier spike is
not reproduced. All 320 outputs, stop states, and normalized schedules match
the original Host arm. The diagnostic completes all requests in 54.394 s,
with 108694 output tokens and 864346 scheduled positions. It changes no
recovery decisions and is not a policy benefit estimate. The call includes
possible deferred STORE startup, but excludes other connector operations,
including preemption handling. Following the declared one-run limit, we stop
this diagnostic: the old spike remains unexplained, and does not justify an
online state predictor or a claim that all outside-call time is computation.

\subsection{Related work and missing strong comparisons}
CacheOPT profiles reconstruction and swap costs as functions of sequence
length, and combines that selection with other memory and scheduling
mechanisms~\cite{cacheopt}. Stream2LLM studies streaming prefill and uses
hardware profiles to choose recomputation or swap at preemption~\cite{stream2llm}.
Our initial algorithm audit used their v1 texts; a targeted check of CacheOPT
v2 and Stream2LLM v3 retains the same static, roundtrip selection boundary.
The code audit uses Stream2LLM artifact commit \texttt{5e2763d51368390ec747945436c7612236258366}.
The inspected artifact compares static cost functions at preemption, fits
quadratic profiles, and combines swap-in and swap-out samples. Its scheduler
argument comes from the current allocation demand, not directly from the
victim's complete known prefix. We do not silently reinterpret it as our
resume-time selector.

Only a partial, selector-only adaptation has been specified. A fair adaptation
must measure remaining native Host and Recompute paths on this device, include
partial host coverage, and exclude already-paid STORE from incremental
comparison. Copying a round-trip curve, halving it, or subtracting an estimated
STORE time does not produce such a measurement. No complete CacheOPT or
Stream2LLM reproduction or sufficiently covered calibrated strong baseline has
been executed here. L1792 must not be presented as that missing comparator.

Prior characterization already reports that concurrency and KV capacity can
change cache/recomputation behavior~\cite{unboxkv}. Cake coordinates concurrent
loading and computation with scheduling changes~\cite{cake}; CacheGen adapts
compressed KV streaming and text-based rebuilding to bandwidth~\cite{cachegen}.
LUMEN considers load-aware failure recovery and can trade checkpoint locality
for recomputation elsewhere~\cite{lumen}. These establish related concerns
with resources, state reuse, and interference. Our fixed-target native choice
does not justify claiming the first resource-aware recovery method, and these
systems are not equivalent binary baselines under our fixed mechanisms.

\subsection{What the present evidence supports}
The matched interventions support a concrete account of token-budget
displacement and changed re-preemption under this runtime. They do not show
that shorter next-output latency is always harmful, that either fixed action
dominates, or that current state predicts an action's net value after
controlling prefix and coverage. Whole-policy runs are not same-state
counterfactuals for each later event. Request repetitions, shared next outputs,
and two runs per policy are not independent statistical replications.

Numerical correctness beyond continuity, full instrumentation cost, a strong
neighbor comparison, an independent workload with actual legal choices, and
the same-work timing root cause remain unresolved. The adapter span in the G
group is 17.4--31.2 ms per arm, but it excludes other hooks, logging, and
observers; it is not total overhead. No externally justified SLO contract was
frozen, so we report raw completion and stalls rather than inventing goodput
thresholds after observing results. These limitations prevent a quality-adjusted
service-benefit or submission-readiness claim. The useful present result is
the demonstrated need to follow a recovery action through subsequent
preemption, peer opportunities, and complete natural-output service.

\paragraph{Evidence and reproducibility.}
All claims above refer to completed local records. The initial event groups are
\path{oct07_one_event_r01} and \path{oct07_e223_r01}; the policy groups are
\path{oct07_budget_r01} and \path{oct07_headroom_r01}; the timing controls are
\path{oct07_telemetry_r01} and \path{oct07_cpu_tail_r01}. The fixed-count E246
pairs are \path{oct08_e246_once53005_probe_r01} and
\path{oct08_e246_once53005_reverse_r01}. Inputs, source hashes,
raw outputs, native schedules, failures, and analysis receipts are retained
under the experiment directory. The canonical \path{RESULTS.md} links the
current summaries. Artifact identities and quality rules are frozen before
their corresponding comparisons; no unrun successor contributes a result.

\begin{thebibliography}{9}
\bibitem{cacheopt}
H. Shen et al.
\newblock Mitigating KV Cache Competition to Enhance User Experience in LLM Inference.
\newblock arXiv:2503.13773v2, March 24, 2025. Relevant choice and coverage sections checked against v2. \url{https://arxiv.org/abs/2503.13773v2}.

\bibitem{stream2llm}
R. Bachkaniwala et al.
\newblock Stream2LLM: Overlap Context Streaming and Prefill for Reduced Time-to-First-Token.
\newblock arXiv:2604.16395v3, May 17, 2026; the author record reports acceptance
at MLSys 2026. Relevant choice and coverage sections checked against v3.
\url{https://arxiv.org/abs/2604.16395v3}.
Official artifact: \url{https://github.com/rajveerb/stream2llm/tree/mlsys_artifact}.

\bibitem{unboxkv}
\newblock Characterizing the Behavior and Impact of KV Caching on Transformer
Inferences under Concurrency.
\newblock Author-hosted full text, 2025.
\url{https://jcernuda.com/publications/kv-caching}.

\bibitem{cake}
S. Jin et al.
\newblock Compute Or Load KV Cache? Why not Both?
\newblock arXiv:2410.03065v2, 2025.
\url{https://arxiv.org/html/2410.03065v2}.

\bibitem{cachegen}
Y. Liu et al.
\newblock CacheGen: KV Cache Compression and Streaming for Fast Language Model Serving.
\newblock ACM SIGCOMM, 2024; inspected arXiv:2310.07240v4.
\url{https://arxiv.org/html/2310.07240v4}.

\bibitem{lumen}
Z. Cao et al.
\newblock LUMEN: Coordinated Failure Recovery for Distributed LLM Serving.
\newblock arXiv:2606.17787v1, 2026.
\url{https://arxiv.org/html/2606.17787v1}.

\bibitem{gsm8k}
OpenAI.
\newblock Grade School Math: public test data, selected indices 128--191.
\url{https://github.com/openai/grade-school-math/blob/master/grade_school_math/data/test.jsonl}.

\bibitem{longbench}
LongBench official dataset and evaluation sources.
\newblock Selected complete \texttt{multifieldqa\_en} and \texttt{multi\_news}
records; data revision \texttt{5e628be450b7e67fb7ae6e201bd6d8f7056f7672}.
\url{https://huggingface.co/datasets/zai-org/LongBench/tree/5e628be450b7e67fb7ae6e201bd6d8f7056f7672}.
\end{thebibliography}
\end{document}
